% concatenated from op3_robust_ladders_submission_positioned.tex
\documentclass[11pt]{article}
\usepackage[T1]{fontenc}
\usepackage{amsmath,amssymb,amsthm,mathtools}
\usepackage[margin=1.1in]{geometry}
\usepackage{microtype}
\usepackage[hidelinks]{hyperref}

\theoremstyle{plain}
\newtheorem{theorem}{Theorem}
\newtheorem{lemma}[theorem]{Lemma}
\newtheorem{proposition}[theorem]{Proposition}
\newtheorem{corollary}[theorem]{Corollary}
\theoremstyle{definition}
\newtheorem{definition}[theorem]{Definition}
\newtheorem{remark}[theorem]{Remark}
\numberwithin{theorem}{section}

\newcommand{\R}{\mathbb{R}}
\newcommand{\Z}{\mathbb{Z}}
\newcommand{\F}{\mathbb{F}}
\newcommand{\inner}[2]{\langle #1,#2\rangle}
\newcommand{\norm}[1]{\lVert #1\rVert}
\newcommand{\abs}[1]{\lvert #1\rvert}
\newcommand{\half}{\tfrac12}
\newcommand{\twr}{\operatorname{twr}}
\newcommand{\ldr}{\operatorname{ldr}}
\newcommand{\ldrs}{\operatorname{ldr}^{*}}
\newcommand{\reg}{\operatorname{reg}}
\newcommand{\codim}{\operatorname{codim}}
\newcommand{\Xii}{\Xi}
\newcommand{\Stable}{\mathcal{S}}
\newcommand{\peridx}{\operatorname{per}}
\newcommand{\econst}{\mathrm e}

\title{Robust order properties in the HLM tower construction for arithmetic regularity over \texorpdfstring{$\F_2^{\,n}$}{F2^n}}
\author{Dean Menezes\thanks{The University of Texas at Austin, Austin, Texas 78712, USA.\\
Email: \href{mailto:dean.menezes@utexas.edu}{dean.menezes@utexas.edu}.}}
\date{}
\hypersetup{
  pdftitle={Robust order properties in the HLM tower construction for arithmetic regularity over F2 vector spaces},
  pdfauthor={Dean Menezes},
  pdfsubject={Robust order properties in arithmetic regularity constructions},
  pdfkeywords={arithmetic regularity, order property, half-graph, stability, edit distance, VC dimension}
}
\emergencystretch=2em
\begin{document}
\maketitle

\begin{abstract}
Green's arithmetic regularity lemma over $G=\F_2^{\,n}$ has tower-type
lower bounds.  We show that the order property in the
Hosseini--Lovett--Moshkovitz--Shapira construction is robust under
adversarial edits.  Write
$f=s^{-1}\sum_{i=1}^{s}\mathbf 1_{A_i}$, where
$s=\lfloor1/(16\epsilon)\rfloor$, and let $d_s$ be the top block
dimension.  A ladder of length $r$ in a set $A$ is a pattern
$a_p+b_q\in A$ exactly when $p\le q$.  For every HLM steering system
satisfying the construction's spanning conditions, changing $A_s$ on at
most $\abs G/16$ points leaves a ladder of length $d_s/8$.  The proof
views the $2^{d_s}$ row traces as a binary code of length $8d_s$ and
relative distance greater than $1/4$, then applies Sauer's lemma.  We also
choose the top steering map so that a superlevel set of $f$ retains a
ladder of length at least $d_s/(8\sqrt{s})\ge\twr(s-2)$ after every edit on
at most $\epsilon\abs G$ points.  Finally, robust ladders obstruct
low-index periodic approximation, and sufficiently costly Green
regularity forces positive edit distance from every fixed stable class.
\end{abstract}

\medskip
\noindent\textbf{2020 Mathematics Subject Classification.}
Primary 11B30; Secondary 03C45, 05C35.

\noindent\textbf{Keywords.}
Arithmetic regularity lemma; order property; half-graph; stable arithmetic
regularity; tower lower bound; edit distance; VC dimension; finite vector spaces.

\section{Introduction}\label{sec:intro}

Green's arithmetic regularity lemma says that a function on
$\F_2^{\,n}$ becomes Fourier-uniform on almost all cosets of a subspace
whose codimension depends only on the accuracy parameter~\cite{G}.  The
best general bound has tower growth.  Hosseini, Lovett, Moshkovitz, and
Shapira (HLM) proved that this growth is necessary: for
$s=\lfloor1/(16\epsilon)\rfloor$ they constructed an average
\[
  f=\frac1s\sum_{i=1}^{s}\mathbf 1_{A_i}
\]
for which the only $\epsilon$-regular subspace is the zero subspace, and
the ambient dimension is at least $\twr(s-1)$~\cite{HLM}.

Stable arithmetic regularity suggests a combinatorial explanation for
this cost.  Terry and Wolf proved that, for each fixed $k$, a $k$-stable
subset of $\F_p^{\,n}$ has a polynomial-codimension regularity theorem
with no exceptional cosets~\cite{TW}.  Thus a tower-forcing family cannot
remain $k$-stable for any fixed $k$.  They asked for an explicit,
quantitative order-property configuration in the Green--HLM examples.

We answer that question and prove a stronger robustness statement.  The
top HLM level set $A_s$ contains a ladder of length $d_s$, where $d_s$ is
the top block dimension and is a tower in $1/\epsilon$.  More
significantly, changing a fixed positive proportion of the set does not
remove all long ladders: every edit of density at most $1/16$ leaves one
of length $d_s/8$.  This conclusion uses all rows of the construction,
not one protected configuration.  Their traces form an error-correcting
code, and Sauer's lemma prevents a stable approximation from supporting
enough distinct traces.

The tower lower bound concerns the averaged function $f$, rather than one
level predicate.  We therefore also choose the top steering map so that a
binary superlevel set of $f$ has edit-robust ladder length at least
$d_s/(8\sqrt{s})$.  This quantity is still at least $\twr(s-2)$.  The
robust order property is consequently visible both in the top component
and in a threshold of the actual HLM witness.

We now state the terminology precisely.

\begin{definition}[ladder, order property, and stability]\label{def:op}
A set $A\subseteq G$ has the \emph{$k$-order property} if there are
$a_1,\dots,a_k,b_1,\dots,b_k\in G$ such that
\[
  a_p+b_q\in A\quad\Longleftrightarrow\quad p\le q
  \qquad(1\le p,q\le k).
\]
Such a configuration is a \emph{ladder of length $k$}.  Equivalently, in
the bipartite relation
\[
  R_A(x,y)\quad\Longleftrightarrow\quad x+y\in A
\]
on two labelled copies of $G$, these vertices induce a half-graph of
height $k$.  The set $A$ is \emph{$k$-stable} if it contains no ladder of
length $k$.

The \emph{ladder index} $\ldr(A)$ is the largest $k$ for which $A$
contains a ladder of length $k$; set $\ldr(\emptyset)=0$.  Thus $A$ is
$k$-stable exactly when $\ldr(A)\le k-1$.
\end{definition}

\begin{remark}\label{rem:ladder-bookkeeping}
The two vertex families in a ladder are automatically injective, because
their membership rows and columns are distinct.  Hence
$\ldr(A)\le\abs G$.  Restricting a ladder to its first $k-1$ indices gives
a ladder of length $k-1$, so $k$-stability implies $k'$-stability for every
$k'\ge k$.
\end{remark}

\begin{definition}[robust ladder index and edit distance]\label{def:robust-index}
For $\eta\in[0,1]$, define
\[
  \ldrs_\eta(A):=
  \min_{\,B\subseteq G,\ \abs B\le\eta\abs G}
  \ldr(A\mathbin\triangle B).
\]
Thus $\ldrs_\eta(A)$ is the largest ladder length that survives every
symmetric-difference edit on at most $\eta\abs G$ points.

Let $\Stable_k(G)$ be the family of $k$-stable subsets of $G$.  For every
nonempty family $\mathcal C\subseteq\mathcal P(G)$, put
\[
  \operatorname{dist}_G(A,\mathcal C)
  :=\min_{C\in\mathcal C}\frac{\abs{A\mathbin\triangle C}}{\abs G}.
\]
\end{definition}

\begin{proposition}[edit-distance interpretation]\label{prop:distance}
For every $A\subseteq G$, every $k\ge1$, and every $\eta\in[0,1]$,
\[
  \ldrs_\eta(A)\ge k
  \quad\Longleftrightarrow\quad
  \operatorname{dist}_G\bigl(A,\Stable_k(G)\bigr)>\eta.
\]
\end{proposition}

\begin{proof}
The inequality $\ldrs_\eta(A)\ge k$ says that every set $C$ with
$\abs{A\mathbin\triangle C}\le\eta\abs G$ contains a ladder of length
$k$.  Equivalently, no $k$-stable set lies within edit distance $\eta$ of
$A$.
\end{proof}

\paragraph{Results.}
Let $d=d_s$.  The four parts of the paper are as follows.
\begin{enumerate}
\item Theorem~\ref{thm:levelset} gives an explicit ladder of length $d$ in
  $A_s$.  Here
  \[
    d\ge\twr(s-1)
    \qquad\text{and}\qquad
    d=(1-o(1))\log_2\abs G.
  \]
  This answers the question posed by Terry and Wolf.
\item Theorem~\ref{thm:robust} proves the stronger estimate
  \[
    \operatorname{dist}_G\bigl(A_s,\Stable_{d/8}(G)\bigr)
    >\frac18\left(1-\left(\frac{\econst}{4}\right)^{d/8}\right)
    >\frac1{16}.
  \]
  Hence every edit of density at most $1/16$ leaves a ladder of length
  $d/8$.  In particular, the robust ladder length is a tower, not merely
  logarithmic in $1/\epsilon$.
\item Theorem~\ref{thm:witness} chooses the top steering map so that a
  superlevel set $S=\{f\ge\theta\}$ satisfies
  \[
    \ldrs_\epsilon(S)\ge\frac{d}{8\sqrt{s}}\ge\twr(s-2).
  \]
  Thus a binary threshold of the averaged witness remains tower-unstable
  after every $\epsilon$-sparse edit.
\item Section~\ref{sec:consequences} gives two general consequences.
  Robust ladders obstruct approximation by low-index unions of cosets.
  Conversely, stable coset-approximation theorems imply that a sufficiently
  large Green regularity index forces positive edit distance from every
  fixed stable class.
\end{enumerate}

The explicit ladder uses \emph{support separation}: one family occupies the
blocks preceding a level, and the other occupies that level.  The nonlinear
steering map then receives an argument from only one family, so membership
reduces to a bilinear pairing.  The robust result has a different source.
The HLM spanning condition makes the row traces a binary code of relative
distance greater than $1/4$; a stable relation has few traces by Sauer's
lemma.  The disagreement between those two counts forces many edits.

\paragraph{Related notions.}
Several quantitative notions make stability insensitive to a negligible
error or to a small density of witnesses.  Terry and Wolf considered
quantitative relaxations in stable arithmetic regularity~\cite[\S5]{TW};
Coregliano and Malliaris introduced an approximate notion in countable
Ramsey theory~\cite[Definition~3.5]{CM}; Martin-Pizarro, Palac\'in, and Wolf
defined almost sure stability by requiring a negligible density of
half-graphs~\cite{MPPW}; and Gir\'on developed regularity theory for
relations stable modulo an ideal~\cite{Giron}.  Gladkova has recently
obtained other arithmetic-regularity lower bounds from the Green--HLM
constructions~\cite{Gladkova}.

Our parameter instead measures normalized adversarial edit distance from
the class of $k$-stable \emph{Cayley} relations.  An edit preserves the
Cayley form: toggling $z\in A$ toggles exactly the pairs $(x,y)$ with
$x+y=z$.  Proposition~\ref{prop:distance} gives the exact edit-distance
interpretation.  Up to reversing the order convention, the unedited
ladder index is the threshold dimension of the row family in
learning-theoretic terminology~\cite{GGKM}.  We do not claim that edit
distance characterizes Green's regularity index; the general implication
proved in Section~\ref{sec:consequences} is one-way.

\paragraph{Conventions.}
The group $G=\F_2^{\,n}$ is written additively, with
$\inner{u}{v}=\sum_j u_jv_j\in\F_2$.  For a statement $P$, write
$[P]\in\{0,1\}$ for its indicator.  We set $\twr(0)=1$ and
$\twr(r+1)=2^{\twr(r)}$.  All asymptotic notation refers to
$\epsilon\to0$.  Throughout the paper we assume $\epsilon\le1/256$ and put
\[
  s:=\left\lfloor\frac1{16\epsilon}\right\rfloor\ge16.
\]
The definition gives
\begin{equation}\label{eq:epsilon-s}
  \frac1{16(s+1)}<\epsilon\le\frac1{16s}.
\end{equation}

\section{The HLM construction}\label{sec:hlm}

We first fix the normalization of Green regularity used throughout the
paper.

\begin{definition}[$\epsilon$-regular subspace and regularity index]\label{def:reg}
Let $H\le G$.  For $x,t\in G$, put
\[
  \widehat f_{x+H}(t):=\frac1{\abs H}
  \sum_{h\in H}f(x+h)(-1)^{\inner{t}{h}}.
\]
The coset $x+H$ is \emph{$\epsilon$-uniform} for $f$ if
$\abs{\widehat f_{x+H}(t)}\le\epsilon$ for every $t\notin H^\perp$, where
\[
  H^\perp:=\{y\in G:\inner{y}{h}=0\text{ for every }h\in H\}.
\]
The subspace $H$ is \emph{$\epsilon$-regular} for $f$ if all but at most
$\epsilon\abs G$ points $x\in G$ lie in $\epsilon$-uniform cosets.  Finally,
\[
  \reg_\epsilon(f):=
  \min\{2^{\codim H}:H\le G\text{ is $\epsilon$-regular for }f\}.
\]
For a set $A\subseteq G$, we apply the definition to $f=\mathbf 1_A$.
\end{definition}

We now record the HLM construction and prove the regularity statement used
later.

Define block dimensions $d_1=1$ and
\[
  d_{i+1}=2^{D_i}\ \ (i=1,2),\qquad d_{i+1}=2^{D_i-3}\ \ (i\ge3),
\]
where $D_i:=\sum_{j\le i}d_j$ is the running total ($D_0=0$).  Thus
\[
  d_1,\dots,d_5=1,2,8,2^8,2^{264},
  \qquad D_1,D_2,D_3=1,3,11.
\]
The sequence displayed in~\cite[\S2.1]{HLM} has these values.  The
piecewise formula on the same page places $i=3$ in the first branch, which
would instead give $d_4=2^{11}$.  The listed sequence and the identity
$2^{D_{i-1}}=8d_i$ used from level $4$ onward both give
$d_4=2^{D_3-3}=2^8$; we follow those two consistent pieces of data.  The
choice of this initial value does not affect any asymptotic conclusion.
Coordinates of $\F_2^{\,n}$ split into $s$ consecutive blocks,
$x=(x^1,\dots,x^s)$ with $x^i\in\F_2^{d_i}$; write
$x^{<i}:=(x^1,\dots,x^{i-1})\in\F_2^{D_{i-1}}$ for the earlier blocks. Set
$n:=D_s$ and $N:=2^{D_{s-1}}=8d_s$; the last equality holds because $s>3$, so
$d_s=2^{D_{s-1}-3}$.

For each $i$, HLM choose an injective steering map
\[
  \xi_i\colon\F_2^{D_{i-1}}\longrightarrow
  \F_2^{d_i}\setminus\{0\},
  \qquad \Xii_i:=\xi_i(\F_2^{D_{i-1}}).
\]
We use the following properties.
\begin{itemize}
\item[(a)] $\xi_i(u)\ne0$ for every $u$.
\item[(b)] $\Xii_i$ spans $\F_2^{d_i}$.
\item[(c)] $\xi_i$ is injective.
\item[(d)] $\xi_i(x^{<i})$ depends only on the blocks preceding $i$.
\item[(e)] If $i>3$, every subset of $\Xii_i$ containing at least
  $3\abs{\Xii_i}/4$ vectors spans $\F_2^{d_i}$.
\end{itemize}
For $i\le3$, the image $\Xii_i$ is a basis; for $i>3$, it is the set of
$8d_i$ vectors supplied by~\cite[Claim~2.1]{HLM}.  Thus (a)--(d) hold at
every level, while (e) is asserted only where the dimensions permit it.  We
call a family of steering maps with these properties an \emph{admissible HLM
realization}.  The level-set ladder uses (a)--(d), and the robust argument
uses (e) only at the top level $s>3$.

Define the level sets and the averaged witness by
\[
  A_i:=\bigl\{x\in\F_2^{\,n}:\inner{x^i}{\xi_i(x^{<i})}=0\bigr\},
  \qquad B_i:=\mathbf 1_{A_i},\qquad
  f:=\frac1s\sum_{i=1}^s B_i\;\in[0,1].
\]
The next proposition repeats the HLM regularity argument, treating the first
three levels separately.

\begin{proposition}[regularity of every admissible realization]
\label{prop:hlm-regularity}
Suppose that the steering maps satisfy (a)--(d), that $\Xii_i$ is a basis for
$i\le3$, and that (e) holds for $i>3$. Then the only
$\epsilon$-regular subspace for $f$ is $\{0\}$. Consequently
\[
  \reg_\epsilon(f)=2^n\ge\twr(s).
\]
\end{proposition}

\begin{proof}
The zero subspace is regular because it has no nontrivial frequencies.  Let
$H\ne\{0\}$, and choose the least $i$ for which some $v\in H$ has
$v^i\ne0$.  Every element of $H$ then has zero blocks before $i$.

\smallskip\noindent\emph{Good prefixes.}
For $g\in G$, let $\gamma_g$ be the vector whose only nonzero block is block
$i$, where it equals $\xi_i(g^{<i})$. Call $g$ good if
$\gamma_g\notin H^\perp$. If $i\le3$, then $2^{D_{i-1}}=d_i$, and $\xi_i$ maps the prefixes
bijectively onto the basis $\Xii_i$.  The nonzero functional
$z\mapsto\inner{z}{v^i}$ is nonzero on at least one basis vector.  Hence at
least a $1/d_i\ge1/8$ fraction of the prefixes are good.  If $i>3$, the bad
steering vectors lie in the hyperplane
$(v^i)^\perp$. Property (e) therefore implies that fewer than three quarters
of the prefixes are bad. In either case at least a fraction
\[
  \rho:=\frac18
\]
of the vectors $g$ are good; the blocks after the prefix are unrestricted.

\smallskip\noindent\emph{Averaging one Fourier coefficient.}
Let $V_{>i}$ be the subspace supported on the blocks after $i$.  Fix a good
$g$, put
\[
  \sigma_g:=(-1)^{\inner{g^i}{\xi_i(g^{<i})}},
\]
and define
\[
  X_w:=\sigma_g\,\widehat f_{H+g+w}(\gamma_g)\qquad(w\in V_{>i}).
\]
The factor $\sigma_g$ chooses the sign of the Fourier coefficient; in
particular, $\abs{X_w}=\abs{\widehat f_{H+g+w}(\gamma_g)}$. Averaging over $w$
gives
\begin{equation}\label{eq:hlm-average}
  \mathbb E_{w\in V_{>i}}X_w=\frac1{2s}.
\end{equation}
Indeed, the terms $B_j$ with $j<i$ are constant on $H+g+w$ and hence have
zero Fourier coefficient at the nontrivial frequency $\gamma_g$. For $j=i$,
write $\chi(h)=(-1)^{\inner{\gamma_g}{h}}$. On the coset $H+g+w$ we have
\[
  B_i(g+w+h)=\frac12\bigl(1+\sigma_g\chi(h)\bigr),
\]
so multiplication by $\sigma_g\chi(h)$ and averaging over $h$ gives $1/2$.
For $j>i$, average first over the fresh block $w^j$; since $\xi_j$ is nonzero
and does not depend on $w^j$, that average is $1/2$, and the remaining
nontrivial character on $H$ has mean zero. This proves
\eqref{eq:hlm-average}.

\smallskip\noindent\emph{Many nonuniform cosets.}
A nontrivial character divides $H$ into two equal halves.  Since $0\le f\le1$,
we have $\abs{X_w}\le1/2$.  Let $q$ be the fraction of $w\in V_{>i}$ for
which $X_w>\epsilon$. Each such $w$ gives a nonuniform coset because
$\abs{\widehat f_{H+g+w}(\gamma_g)}>\epsilon$. From
\eqref{eq:hlm-average},
\[
  \frac1{2s}\le \frac q2+(1-q)\epsilon,
\]
so, by \eqref{eq:epsilon-s},
\[
  q\ge\frac{1/(2s)-\epsilon}{1/2-\epsilon}
  \ge\frac7{8s}.
\]
Choose $g$ uniformly from $G$ and $w$ uniformly from $V_{>i}$.  At least a
$\rho$ fraction of the choices of $g$ are good; conditional on each good
$g$, at least a $q$ fraction of the choices of $w$ give a nonuniform coset
$H+g+w$.  Therefore
\[
  \Pr_{g,w}(H+g+w\text{ is nonuniform})
  \ge\rho q\ge\frac7{64s}>\frac1{16s}\ge\epsilon.
\]
The sum $g+w$ is uniform on $G$, so the left side is exactly the fraction of
points that lie in nonuniform $H$-cosets.  Hence $H$ is not
$\epsilon$-regular.

It follows that $\reg_\epsilon(f)=2^n$. Since
$n\ge d_s\ge\twr(s-1)$ by Lemma~\ref{lem:blockgrowth},
$2^n\ge\twr(s)$.
\end{proof}

We also need the following elementary estimates for the block dimensions.

\begin{lemma}\label{lem:blockgrowth}
For every $i\ge1$ one has $d_i\ge\twr(i-1)$, and $\twr(m)\ge2^{m}$ for every
$m\ge1$. Consequently $d_s\ge\twr(s-1)$; moreover $n=d_s+\log_2 d_s+3$, so
$d_s\ge\half\log_2\abs{G}$, and $d_s=(1-o(1))\log_2\abs{G}$ as $\epsilon\to0$.
\end{lemma}

\begin{proof}
The first four cases are
\[
  (d_1,d_2,d_3,d_4)=(1,2,8,2^8)
  \quad\text{and}\quad
  (\twr(0),\twr(1),\twr(2),\twr(3))=(1,2,4,16).
\]
For $i\ge4$, since $D_{i-1}\ge D_3=11>3$
we have $D_i-3=d_i+(D_{i-1}-3)\ge d_i$, so
\[
  d_{i+1}=2^{D_i-3}\ge 2^{d_i}\ge 2^{\twr(i-1)}=\twr(i).
\]
The bound $\twr(m)\ge2^m$ also goes by induction: $\twr(1)=2$, and
$\twr(m+1)=2^{\twr(m)}\ge2^{2^m}\ge2^{m+1}$.

For the final claims, the standing assumption gives $s>3$, so
$d_s=2^{D_{s-1}-3}$, that is, $D_{s-1}=\log_2 d_s+3$ and $N=8d_s$; hence
$n=D_s=d_s+D_{s-1}=d_s+\log_2 d_s+3$. Since $d_s\ge 8$ we have
$d_s\ge\log_2 d_s+3$, so $2d_s\ge n$ and $d_s\ge\half\log_2\abs{G}$. Finally
$d_s/n=1-(\log_2 d_s+3)/n$, and $d_s\ge\twr(s-1)\to\infty$ as $\epsilon\to0$,
so $d_s/n\to1$.
\end{proof}



\section{An explicit ladder in a level set}\label{sec:levelset}

The predicate defining $A_i$ is the bilinear expression
$\inner{x^i}{\xi_i(x^{<i})}$, except that the second argument is chosen by the
nonlinear steering map $\xi_i$. We place the two index families in disjoint
blocks. Then $\xi_i$ receives an argument from only one family, and membership
reduces to an ordinary bilinear pairing.

\begin{theorem}[level-set ladder; for every admissible realization]\label{thm:levelset}
For every $1\le i\le s$ and every $k\le d_i$, explicit vectors witness
the $k$-order property in $A_i$. In particular, $A_s$ has the $d_s$-order
property, where
\[
  d_s\ \ge\ \twr(s-1)=\twr\bigl(\lfloor 1/(16\epsilon)\rfloor-1\bigr)
  \qquad\text{and}\qquad d_s\ge\half\log_2\abs{G},
\]
and $d_s=(1-o(1))\log_2\abs{G}$ as $\epsilon\to0$.  In particular, for each
fixed $k$, the set $A_s$ fails to be $k$-stable once $\epsilon$ is
sufficiently small.
\end{theorem}

\begin{proof}
Fix $i$ and $k\le d_i$.

\emph{Columns.} By property (b), $\Xii_i$ spans $\F_2^{d_i}$, so it contains
$d_i$ linearly independent vectors; choose linearly independent
$v_1,\dots,v_k\in\Xii_i$.  Choose their preimages $u_\ell$, so that
$\xi_i(u_\ell)=v_\ell$.  The vectors $u_\ell$ are distinct because $\xi_i$
is injective.

\emph{Rows.} Extend $v_1,\dots,v_k$ to a basis $v_1,\dots,v_{d_i}$ of
$\F_2^{d_i}$. Since the standard form is nondegenerate, the map
$y\mapsto\inner{y}{\cdot}$ is an isomorphism onto the dual space, so the basis
has a dual basis $v_1^{*},\dots,v_{d_i}^{*}$, characterized by
$\inner{v_j^{*}}{v_\ell}=[j=\ell]$. For $1\le m\le k$ set
$w_m:=\sum_{\ell=m+1}^{k}v_\ell^{*}$ (an empty sum for $m=k$, so $w_k=0$);
then for $1\le\ell\le k$,
\[
  \inner{w_m}{v_\ell}=\sum_{j=m+1}^{k}[j=\ell]=[\ell>m].
\]

\emph{Separation.} In block coordinates set $a_\ell:=(u_\ell,0)$, carrying
$u_\ell$ in the blocks $<i$ and zero elsewhere, and $b_m:=(0,w_m,0)$, carrying
$w_m$ in block $i$ and zero elsewhere. Because addition is coordinatewise and
the blocks are disjoint, the sum $x:=a_\ell+b_m$ has $x^{<i}=u_\ell$ and
$x^i=w_m$; the blocks after the $i$th are zero and do not enter the definition
of $A_i$. Thus $\xi_i$ is evaluated at $u_\ell$, which depends only on $\ell$; no
additivity of $\xi_i$ is needed. Therefore
\[
  \inner{x^i}{\xi_i(x^{<i})}=\inner{w_m}{\xi_i(u_\ell)}
  =\inner{w_m}{v_\ell}=[\ell>m].
\]

\emph{Ladder.} Consequently $a_\ell+b_m\in A_i\iff\inner{w_m}{v_\ell}=0\iff
\ell\le m$, which is the $k$-order property of Definition~\ref{def:op} with the
$a$'s indexed by $\ell$ and the $b$'s by $m$. The $a_\ell$ are distinct because
the $u_\ell$ are, and the $w_m$ are distinct because they realize the distinct
patterns $([\ell>m])_{\ell=1}^{k}$.

\emph{Length.} Take $i=s$ and $k=d_s$; the stated bounds on $d_s$ are
Lemma~\ref{lem:blockgrowth}.
\end{proof}

\begin{remark}[a hand instance]\label{rem:hand}
With $v_\ell=e_\ell$ and $d=3$, the construction gives
$w_1=(0,1,1)$, $w_2=(0,0,1)$, and $w_3=(0,0,0)$.  The $3\times3$ membership
table of $a_\ell+b_m\in A_i$ is
exactly the half-graph $[\ell\le m]$. That $w_3=0$ is harmless: probe vectors
need not be nonzero, only the steering values $\xi_i(\cdot)$ are, and those are
the $v_\ell$.
\end{remark}

\begin{remark}[the support-separated maximum]\label{rem:exact}
Consider the relation
\[
  \inner{w}{v}=0
  \qquad\bigl(w\in\F_2^d,\ v\in\F_2^d\setminus\{0\}\bigr).
\]
A ladder $\inner{w_m}{v_\ell}=[\ell>m]$ forces
$v_1,\dots,v_k$ to be linearly independent.  Indeed, suppose
$\sum_{\ell\in S}v_\ell=0$ for a nonempty set $S$.  A singleton $S$ would
give a zero column.  Otherwise let $s_1<s_2$ be the two least elements of
$S$.  Pairing the dependence with $w_{s_1}$ and $w_{s_2}$ gives
\[
  0\equiv\abs S-1\pmod2,
  \qquad
  0\equiv\abs S-2\pmod2,
\]
a contradiction.  Hence this relation has ladder index exactly $d$, and
Theorem~\ref{thm:levelset} is optimal for the support-separated mechanism.
If zero columns are allowed, the maximum becomes $d+1$: only the first
column can be zero, the remaining columns are independent, and prepending a
zero column to the construction above attains this bound.  We do not
claim that $\ldr(A_s)=d_s$ in the full Cayley relation; we prove only
$\ldr(A_s)\ge d_s$.
\end{remark}

\section{A robust ladder in the top level set}\label{sec:robust}

Write
\[
  W:=\F_2^{D_{s-1}},\qquad U:=\F_2^{d},\qquad d:=d_s,
  \qquad N:=\abs W=8d,
\]
so that $G=W\oplus U$.  For a set $C\subseteq W\times U$ and $u\in U$, its
\emph{row trace} is
\[
  R_C(u):=\{w\in W:(w,u)\in C\}\subseteq W.
\]
The VC dimension of the row family $\{R_C(u):u\in U\}$ is the largest
integer $r$ for which some $r$-element subset of $W$ is shattered.  Shattering
$r$ columns already gives a ladder of length $r$, because the row family
realizes the $r$ suffix patterns $([p\le q])_{q=1}^r$.  We combine this
observation with Sauer's lemma.

\begin{lemma}[Sauer bound for stable row families]\label{lem:trace-bound}
If $C\subseteq W\times U$ has no ladder of length $r$, then
\[
  \abs{\{R_C(u):u\in U\}}
  \le \sum_{j=0}^{r-1}\binom Nj.
\]
Moreover, for integers $1\le q\le N/2$,
\[
  \sum_{j=0}^{q}\binom Nj\le\left(\frac{\econst N}{q}\right)^q.
\]
\end{lemma}

\begin{proof}
If $w_1,\dots,w_r$ are shattered, choose $u_p$ whose trace on these columns
is $([p\le q])_{q=1}^r$.  Then $(0,u_p)+(w_q,0)\in C$ exactly when $p\le q$,
which is a ladder of length $r$.  Hence the row family has VC dimension at
most $r-1$, and Sauer's lemma gives the first inequality~\cite{Sauer}.  For
the second inequality, put $x=q/(N-q)\le1$.  Since $x^j\ge x^q$ for
$0\le j\le q$,
\[
  x^q\sum_{j=0}^{q}\binom Nj
  \le \sum_{j=0}^{q}\binom Nj x^j
  \le (1+x)^N.
\]
Therefore
\[
  \sum_{j=0}^{q}\binom Nj
  \le \left(\frac Nq\right)^q
       \left(\frac N{N-q}\right)^{N-q}
  \le \left(\frac{\econst N}{q}\right)^q,
\]
because $N/(N-q)=1+q/(N-q)$ and $(1+y)^{1/y}\le\econst$ for $y>0$.
\end{proof}

For $A=A_s$, the row traces are not merely numerous.  They form an
error-correcting code.

\begin{theorem}[constant-edit robustness of the top level]\label{thm:robust}
Let $A:=A_s$ come from an admissible HLM realization.  Because $d$ is a
power of $2$ and $s\ge16$, it is divisible by $8$; put $r:=d/8$.  Then
\begin{equation}\label{eq:robust-distance}
  \operatorname{dist}_G\bigl(A,\Stable_r(G)\bigr)
  >\frac18\left(1-\left(\frac{\econst}{4}\right)^r\right)
  >\frac1{16}.
\end{equation}
Consequently, for every $0\le\eta\le1/16$,
\[
  \ldrs_\eta(A)\ge\frac d8\ge\twr(s-2),
\]
and $d/8\ge(1/16)\log_2\abs G$.  More generally, for each fixed $\eta<1/8$, one has
$\ldrs_\eta(A)\ge d/8$ once $s$ is sufficiently large.
\end{theorem}

\begin{proof}
For $u\in U$, let $\rho_u\in\{0,1\}^{W}$ be the indicator of $R_A(u)$:
\[
  \rho_u(w)=[\inner{u}{\xi_s(w)}=0].
\]
If $u\ne v$, put $z=u+v\ne0$.  The two rows differ exactly at the columns
for which $\inner{z}{\xi_s(w)}=1$.  If they differed in at most $N/4$
columns, then at least $3N/4$ steering vectors would lie in the proper
hyperplane $z^\perp$.  Property~\textup{(e)} says that those vectors span
$U$, a contradiction.  Hence, writing $d_H$ for Hamming distance,
\begin{equation}\label{eq:code-distance}
  d_H(\rho_u,\rho_v)>\frac N4\qquad(u\ne v).
\end{equation}
Thus the $2^d$ rows form a binary code of length $N=8d$ and relative
distance greater than $1/4$.

Let $C\subseteq G$ be $r$-stable, and let
$M:=\abs{\{R_C(u):u\in U\}}$ be the number of its distinct row traces.
Lemma~\ref{lem:trace-bound}, with $N=8d$ and $d=8r$, gives
\begin{equation}\label{eq:stable-traces}
  M\le\sum_{j=0}^{r-1}\binom Nj
  \le(64\econst)^r
  =2^d\left(\frac{\econst}{4}\right)^r.
\end{equation}
Associate each original row $\rho_u$ with its edited trace $R_C(u)$.  A fixed
edited trace can lie within Hamming distance $N/8$ of at most one codeword:
two such codewords would be at distance at most $N/4$, contrary to
\eqref{eq:code-distance}.  Since $C$ has only $M$ distinct traces, at most
$M$ original rows are changed in at most $N/8$ positions.  Every remaining
row requires more than $N/8$ edits.  Hence
\[
  \frac{\abs{A\mathbin\triangle C}}{\abs G}
  >\frac{(2^d-M)N/8}{2^dN}
  =\frac18\left(1-\frac{M}{2^d}\right)
  \ge\frac18\left(1-\left(\frac{\econst}{4}\right)^r\right).
\]
Taking the minimum over $C\in\Stable_r(G)$ proves the first inequality in
\eqref{eq:robust-distance}.  Since $r=d/8\ge3$ and $\econst<3$,
$(\econst/4)^r<(3/4)^3<1/2$, which proves the second.
Proposition~\ref{prop:distance} now gives the robust bound for
$\eta\le1/16$.

Lemma~\ref{lem:blockgrowth} gives
$d/8\ge(1/16)\log_2\abs G$.  For the tower bound, put
$t:=\twr(s-2)$.  Then $d\ge\twr(s-1)=2^t$, while $2^t\ge8t$ for
$t\ge4$.  Hence $d/8\ge t=\twr(s-2)$.  Finally, the right side of
\eqref{eq:robust-distance} tends to $1/8$ as $s\to\infty$, which proves the
last assertion.
\end{proof}

\begin{remark}\label{rem:log-false}
Protecting one prescribed ladder row by row gives only a logarithmic bound.
Theorem~\ref{thm:robust} instead uses all $2^d$ rows.  Their code distance
prevents a sparse edit from merging many traces, and Sauer's lemma then
forces a ladder of linear length.  Thus no
$O(\log(1/\epsilon))$ upper bound holds for
$\ldrs_\epsilon(A_s)$ at the HLM scale.
\end{remark}

\section{A robust superlevel set of the averaged witness}\label{sec:witness}

The tower lower bound concerns the average $f$, not one level predicate.  We
now extract a binary superlevel set of $f$ whose ladders also survive sparse
edits.  For $w\in W$, define
\[
  c(w):=\sum_{j<s}B_j(x)
  \qquad\text{for any $x$ with $x^{<s}=w$}.
\]
This is well defined because the lower predicates do not read block $s$.

For uniform $w\in W$, the variables $B_1(w),\dots,B_{s-1}(w)$ are
independent fair bits.  Conditional on the preceding blocks,
$\xi_j(w^{<j})$ is a fixed nonzero vector, whose pairing with the fresh block
$w^j$ is uniform in $\F_2$.  Hence
$c\sim\operatorname{Binomial}(s-1,1/2)$.

\begin{lemma}[a large independent level fibre]\label{lem:independent-fibre}
Fix admissible lower maps $\xi_1,\dots,\xi_{s-1}$.  There are a value
$c_0\in\{0,\dots,s-1\}$, a set $F_0\subseteq c^{-1}(c_0)$, and an admissible
choice of the top map $\xi_s$ such that, with $m:=\abs{F_0}$,
\[
  16\mid m,\qquad \frac{2d}{\sqrt{s}}\le m\le\frac d2,
\]
and the vectors $\{\xi_s(w):w\in F_0\}$ are linearly independent.
\end{lemma}

\begin{proof}
Choose $c_0$ at a mode of the binomial distribution and put
$F=c^{-1}(c_0)$.  We first bound the modal mass.  Define
\[
  a_r:=4^{-r}\binom{2r}{r}\qquad(r\ge1).
\]
Then $a_1=1/2$ and
\[
  a_r=\left(1-\frac1{2r}\right)a_{r-1}.
\]
Induction gives $a_r\ge1/(2\sqrt r)$.  Indeed, for $r\ge2$ the recurrence
and the induction hypothesis give
\[
  a_r\ge\frac{2r-1}{4r\sqrt{r-1}}\ge\frac1{2\sqrt r};
\]
the last inequality follows, after squaring, from
$r(2r-1)^2-4r^2(r-1)=r>0$.  The modal mass for $2r$ trials is $a_r$, and the
modal mass for $2r+1$ trials is $a_{r+1}$.  Applying these bounds with
$s-1$ trials gives
\[
  2^{-(s-1)}\binom{s-1}{\lfloor(s-1)/2\rfloor}
  \ge\frac1{2\sqrt{s}}.
\]
Consequently,
\begin{equation}\label{eq:fibre-size}
  \abs F\ge\frac{N}{2\sqrt{s}}=\frac{4d}{\sqrt{s}}.
\end{equation}
Let $L:=\min\{\abs F,d/2\}$ and set
\[
  m:=16\left\lfloor\frac{L}{16}\right\rfloor.
\]
Since $s\ge16$, the power of two $d=2^{D_{s-1}-3}$ is divisible by $32$;
thus $d/2$ is divisible by $16$.  If $16\le s\le64$, then \eqref{eq:fibre-size} gives $L=d/2$, hence
$m=d/2\ge2d/\sqrt{s}$.  If $s>64$, then both $\abs F$ and $d/2$ are at
least $4d/\sqrt{s}$, so $L\ge4d/\sqrt{s}$.  Also $L\ge32$; hence
$m\ge L-16\ge L/2\ge2d/\sqrt{s}$.  Choose any
$F_0\subseteq F$ of size $m$.

Now draw the vectors $g_w:=\xi_s(w)$ independently and uniformly from
$U=\F_2^d$.  The probability that some $g_w$ is zero is at most $N2^{-d}$, and the
probability of a collision is at most $\binom{N}{2}2^{-d}$.  If
property~\textup{(e)} fails, some hyperplane contains at least $6d$ of the
$8d$ vectors.  For a fixed hyperplane this event has probability at most
$\econst^{-d}$ by Hoeffding's inequality; a union bound over the fewer than
$2^d$ hyperplanes gives the bound $(2/\econst)^d$.  In addition, revealing the vectors indexed by
$F_0$ in order shows that the probability that they are linearly dependent is
less than
\[
  \sum_{j=0}^{m-1}2^{j-d}<2^{m-d}\le2^{-d/2}.
\]
Since $d\ge\twr(15)>256$, the sum of these four bounds is less than $1$.
Choose a realization outside all four bad events.  It satisfies
\textup{(a)}, \textup{(c)}, and \textup{(e)} at level $s$, and hence also
\textup{(b)}; property~\textup{(d)} is built into the definition.  The
chosen vectors on $F_0$ are linearly independent.
\end{proof}

\begin{theorem}[edit-robust superlevel set]\label{thm:witness}
There exist an admissible realization of the HLM witness $f$, a threshold
$\theta=(c_0+1)/s$, and an integer $m$ satisfying
\[
  16\mid m,\qquad \frac{2d}{\sqrt{s}}\le m\le\frac d2,
\]
such that the set
\[
  S:=\{x\in G:f(x)\ge\theta\}
\]
satisfies
\[
  \ldrs_\epsilon(S)\ge\frac{m}{16}
  \ge\frac{d}{8\sqrt{s}}
  \ge\twr(s-2).
\]
Consequently
\[
  \operatorname{dist}_G\bigl(S,\Stable_{m/16}(G)\bigr)>\epsilon.
\]
The same realization still satisfies
$\reg_\epsilon(f)=2^n\ge\twr(s)$.
\end{theorem}

\begin{proof}
Take $c_0,F_0,m$, and $\xi_s$ from Lemma~\ref{lem:independent-fibre}.  On the
rectangle $F_0\times U$, the lower levels contribute the constant $c_0$.
Therefore
\begin{equation}\label{eq:threshold-rectangle}
  (w,u)\in S
  \quad\Longleftrightarrow\quad
  \inner{u}{\xi_s(w)}=0
  \qquad(w\in F_0,u\in U).
\end{equation}
Since the $m$ vectors $\xi_s(w)$, $w\in F_0$, are independent, the row traces
of $S$ on $F_0$ run through all $2^m$ binary patterns, each exactly
$2^{d-m}$ times.

Fix an arbitrary edit $B\subseteq G$ with $\abs B\le\epsilon\abs G$.
The rectangle $F_0\times U$ has size $m2^d$, whereas $G$ has size
$N2^d$.  Hence the relative edit density on this rectangle is at most
\[
  \delta:=\frac{\epsilon N}{m}.
\]
If $s\le64$, then $m=d/2$, and \eqref{eq:epsilon-s} gives
$\delta\le16\epsilon\le1/s\le1/16$.  If $s>64$, then
$m\ge2d/\sqrt{s}$, so again by \eqref{eq:epsilon-s},
$\delta\le4\epsilon\sqrt{s}\le1/(4\sqrt{s})<1/16$.  Thus
\begin{equation}\label{eq:local-edit}
  \delta\le\frac1{16}.
\end{equation}

The average number of edited positions in $F_0$ per row is at most
$\delta m$.  Markov's inequality therefore leaves at least half of the
$2^d$ rows with at most $2\delta m\le m/8$ edits.  Call these rows good.  Every original
trace has $2^{d-m}$ representatives, so at least $2^{m-1}$ distinct original
traces have a good representative.  Choose one good row for each such trace.
Every chosen original trace lies within Hamming distance $m/8$ of its
edited trace.  Therefore one edited trace can arise from at most
\[
  V:=\sum_{j=0}^{m/8}\binom mj
  \le(8\econst)^{m/8}
\]
of the chosen original traces.  Hence the edited relation has at least
$2^{m-1}/V$ distinct row traces on $F_0$.

Suppose that $S\mathbin\triangle B$ has no ladder of length $r:=m/16$.
Lemma~\ref{lem:trace-bound}, applied to the $m$ columns in $F_0$, bounds the
number of distinct edited traces by
\[
  \sum_{j=0}^{r-1}\binom mj\le(16\econst)^{m/16}.
\]
It follows that
\[
  2^{m-1}
  \le(8\econst)^{m/8}(16\econst)^{m/16}
  =2^m\left(\frac{\econst^3}{64}\right)^{m/16}
  <2^{m-1},
\]
because $m/16\ge1$ and $\econst^3/64<27/64<1/2$.  This contradiction proves
$\ldr(S\mathbin\triangle B)\ge m/16$.  Since $B$ was arbitrary,
$\ldrs_\epsilon(S)\ge m/16$.  The lower bound
$m/16\ge d/(8\sqrt{s})$ is Lemma~\ref{lem:independent-fibre}.
For the tower bound, put $t:=\twr(s-2)$.  Then $d\ge2^t$ and $s\le t$.
Moreover, $2^t\ge8t^{3/2}$ for $t\ge16$.  Indeed, the sequence
$2^t/t^{3/2}$ is increasing there, since the ratio of consecutive terms is
$2(t/(t+1))^{3/2}>1$, and the inequality holds at $t=16$.
Therefore
\[
  \frac{d}{8\sqrt{s}}\ge\frac{2^t}{8\sqrt t}\ge t=\twr(s-2).
\]
Proposition~\ref{prop:distance} gives the distance statement.  Finally, the
selected top map is admissible, so Proposition~\ref{prop:hlm-regularity}
gives the regularity index of $f$.
\end{proof}

\section{Consequences for periodic approximation and Green regularity}\label{sec:consequences}

The robust ladder index does not by itself lower-bound Green's Fourier
regularity index.  It does, however, obstruct a natural stronger form of
regularization: approximation by a Boolean function that is constant on
cosets.

\subsection{Periodic Boolean models}

A set $C\subseteq G$ is \emph{$H$-periodic} if it is a union of cosets of
$H\le G$.
\begin{definition}[periodic approximation index]
For $\eta\in[0,1]$, define
\[
  \peridx_\eta(A):=
  \min\bigl\{[G:H]:\text{some $H$-periodic $C$ satisfies }
  \abs{A\mathbin\triangle C}\le\eta\abs G\bigr\}.
\]
The minimum exists because every subset of $G$ is $\{0\}$-periodic.
\end{definition}

\begin{proposition}[robust ladders obstruct periodic approximation]
\label{prop:periodic}
For every $A\subseteq G$ and every $\eta\in[0,1]$,
\[
  \peridx_\eta(A)\ge\ldrs_\eta(A).
\]
More generally, let $0\le\alpha,\beta\le1$ with $\alpha+\beta\le1$.
Suppose that the union of the exceptional $H$-cosets has size at most
$\beta\abs G$, and that every other $H$-coset has $A$-density at most
$\alpha$ or at least $1-\alpha$.  Then
\[
  [G:H]\ge\ldrs_{\alpha+\beta}(A).
\]
\end{proposition}

\begin{proof}
If $C$ is $H$-periodic, membership in $a+b\in C$ depends only on the two
classes of $a$ and $b$ in $G/H$.  The left vertices of a ladder have distinct
membership rows, hence distinct classes modulo $H$.  Therefore
$\ldr(C)\le[G:H]$.  If $\abs{A\mathbin\triangle C}\le\eta\abs G$, the edit
$A\mathbin\triangle C$ is admissible in the definition of
$\ldrs_\eta(A)$, so
\[
  \ldrs_\eta(A)\le\ldr(C)\le[G:H].
\]
Taking the minimum proves the first assertion.

For the second, on each nonexceptional coset choose the constant value that
agrees with the denser side of $A$; choose either value on exceptional
cosets.  The resulting $H$-periodic set differs from $A$ on at most an
$(\alpha+\beta)$-fraction of $G$.  Apply the first assertion.
\end{proof}

The two robust theorems now give tower-size obstructions to Boolean coset
models:
\begin{equation}\label{eq:periodic-applications}
  \peridx_{1/16}(A_s)\ge\frac d8,
  \qquad
  \peridx_\epsilon(S)\ge\frac{d}{8\sqrt{s}}.
\end{equation}
In particular, any subspace whose cosets approximate either set at the
stated accuracy must itself have tower-size index.

\subsection{Distance from stability and Fourier regularity}

We next combine the robust index with known coset-approximation theorems.
The translate family of a $k$-stable set has VC dimension at most $k-1$:
if it shattered $k$ points, rows realizing the $k$ suffix patterns would
form a ladder of length $k$.  Alon, Fox, and Zhao proved that, in a finite
abelian group of bounded exponent, every set whose translate family
has VC dimension at most $d_0$ can be edited on an $\alpha$-fraction of the
group into a union of cosets of a subgroup of index
$\alpha^{-d_0-o(1)}$~\cite{AFZ}.  Conant later obtained a bound
$\alpha^{-O_k(1)}$ for $k$-stable sets in arbitrary finite groups~\cite{Conant}.

For $k\ge2$ and $0<\alpha<1$, let $Q_k(\alpha)$ be any uniform bound with
the following property.  Every $k$-stable subset of every finite elementary
abelian $2$-group is within normalized edit distance $\alpha$ of a union of
cosets of a subspace of index at most $Q_k(\alpha)$.  The cited results show
that such a bound exists and is polynomial in $1/\alpha$ when $k$ is fixed.

\begin{proposition}[Green regularity cost forces distance from stability]
\label{prop:regularity-distance}
Let $A\subseteq\F_2^{\,n}$, let $k\ge2$, and let
$0<\alpha,\delta<1$ with $\alpha+\delta\le1$.  If
\[
  \operatorname{dist}_G\bigl(A,\Stable_k(G)\bigr)\le\delta,
\]
then
\[
  \reg_{\sqrt{\alpha+\delta}}(\mathbf1_A)\le Q_k(\alpha).
\]
In particular, for $0<\rho<1$,
\[
  \reg_\rho(\mathbf1_A)>Q_k(\rho^2/2)
  \quad\Longrightarrow\quad
  \operatorname{dist}_G\bigl(A,\Stable_k(G)\bigr)>\frac{\rho^2}{2}.
\]
\end{proposition}

\begin{proof}
Choose a $k$-stable set $A'$ with
$\abs{A\mathbin\triangle A'}\le\delta\abs G$.  By the definition of
$Q_k(\alpha)$, there are a subspace $H$ with $[G:H]\le Q_k(\alpha)$ and an
$H$-periodic set $C$ such that
\[
  \abs{A'\mathbin\triangle C}\le\alpha\abs G.
\]
Thus $\abs{A\mathbin\triangle C}\le(\alpha+\delta)\abs G$.  Put
$\gamma:=\alpha+\delta$.  For each coset $x+H$, let $e(x+H)$ be the relative
edit density between $A$ and $C$ on that coset.  The average of $e$ over the
cosets is at most $\gamma$, so all but a $\sqrt\gamma$-fraction of the points
lie in cosets with $e\le\sqrt\gamma$.

On such a coset, $\mathbf1_C$ is constant.  Its Fourier coefficient at every
$t\notin H^\perp$ is therefore zero, while
\[
  \left|\widehat{\mathbf1_A}_{x+H}(t)
        -\widehat{\mathbf1_C}_{x+H}(t)\right|
  \le e(x+H)\le\sqrt\gamma.
\]
Hence $H$ is $\sqrt\gamma$-regular for $\mathbf1_A$, which proves the first
claim.  Set $\alpha=\delta=\rho^2/2$ and take the contrapositive to obtain the
second.
\end{proof}

The implication is one-way.  Distance from every $k$-stable set need not, by
itself, produce large Fourier coefficients on many cosets.  The exact
consequence of a robust ladder is Proposition~\ref{prop:periodic}: it rules
out low-index, nearly Boolean coset models.

\section{Discussion and open questions}\label{sec:remarks}

The two robust theorems settle two questions.  First,
$\ldrs_\epsilon(A_s)$ is not $O(\log(1/\epsilon))$:
Theorem~\ref{thm:robust} gives the lower bound $d_s/8$, even after edits of
fixed density $1/16$.  Second, Theorem~\ref{thm:witness} produces a binary
superlevel set of the averaged witness whose edit-robust ladder length is at
least $d_s/(8\sqrt{s})$.  Both lower bounds are towers in $1/\epsilon$.

Several quantitative questions remain.
\begin{enumerate}
\item Determine the correct order of $\ldrs_\eta(A_s)$ for fixed
  $0<\eta<1/8$.  Is it $\Theta(d_s)$?  The lower bound is linear in $d_s$,
  but ladders in the full Cayley relation need not be support-separated, so
  the argument of Remark~\ref{rem:exact} gives no matching upper bound.
\item Can the factor $\sqrt{s}$ in Theorem~\ref{thm:witness} be removed?
  The loss comes from choosing one level fibre of the binomial variable
  $c(w)$.  A construction that combines several fibres without losing a
  fixed threshold might yield a robust ladder of order $d_s$.
\item Proposition~\ref{prop:regularity-distance} says that expensive Fourier
  regularity forces distance from each fixed stable class.  Find useful
  hypotheses under which a converse holds, or construct sharp counterexamples.
\item Compare adversarial edit distance from $k$-stability with almost sure
  stability and stability modulo an ideal~\cite{CM,MPPW,Giron}.  Our parameter
  asks for a global correction; the other notions control the density or
  negligibility of local witnesses.
\end{enumerate}


\begin{thebibliography}{99}
\small
\raggedright
\bibitem{AFZ}
  N.~Alon, J.~Fox, and Y.~Zhao,
  \emph{Efficient arithmetic regularity and removal lemmas for induced
    bipartite patterns},
  Discrete Anal. (2019), Paper No.~3, 14 pp.,
  doi:10.19086/da.7757; arXiv:1801.04675.
\bibitem{Conant}
  G.~Conant,
  \emph{Quantitative structure of stable sets in arbitrary finite groups},
  Proc.\ Amer.\ Math.\ Soc. \textbf{149} (2021), 4015--4028,
  doi:10.1090/proc/15479; arXiv:2004.02819.
\bibitem{CM}
  L.~N.~Coregliano and M.~Malliaris,
  \emph{Countable Ramsey},
  J. Math. Log., article 2550021 (2025),
  doi:10.1142/S0219061325500217; arXiv:2203.10396.
\bibitem{GGKM}
  B.~Ghazi, N.~Golowich, R.~Kumar, and P.~Manurangsi,
  \emph{Near-tight closure bounds for the Littlestone and threshold
    dimensions},
  Proc.\ Mach.\ Learn.\ Res. \textbf{132} (2021), 686--696;
  arXiv:2007.03668.
\bibitem{Giron}
  M.~Gir\'on,
  \emph{On the regularity of almost stable relations},
  arXiv:2508.00511v2 (2026).
\bibitem{Gladkova}
  V.~Gladkova,
  \emph{A note on lower bounds for arithmetic regularity partitions},
  arXiv:2510.15532 (2025).
\bibitem{G}
  B.~Green,
  \emph{A Szemer\'edi-type regularity lemma in abelian groups, with
    applications},
  Geom.\ Funct.\ Anal. \textbf{15} (2005), no.~2, 340--376,
  doi:10.1007/s00039-005-0509-8.
\bibitem{HLM}
  K.~Hosseini, S.~Lovett, G.~Moshkovitz, and A.~Shapira,
  \emph{An improved lower bound for arithmetic regularity},
  Math.\ Proc.\ Cambridge Philos.\ Soc. \textbf{161} (2016), no.~2, 193--197,
  doi:10.1017/S030500411600013X; arXiv:1405.4409.
\bibitem{MPPW}
  A.~Martin-Pizarro, D.~Palac\'in, and J.~Wolf,
  \emph{Stability, corners, and other two-dimensional shapes},
  Canad. J. Math., published online 26 December 2025,
  doi:10.4153/S0008414X25101971; arXiv:2210.14039.
\bibitem{Sauer}
  N.~Sauer,
  \emph{On the density of families of sets},
  J. Combin. Theory Ser. A \textbf{13} (1972), 145--147,
  doi:10.1016/0097-3165(72)90019-2.
\bibitem{TW}
  C.~Terry and J.~Wolf,
  \emph{Stable arithmetic regularity in the finite-field model},
  Bull.\ London Math.\ Soc. \textbf{51} (2019), no.~1, 70--88,
  doi:10.1112/blms.12211; arXiv:1710.02021.
\end{thebibliography}

\end{document}
